\subsection{模型与环境延续}

第三问沿用\eqref{eq:q2-pde}--\eqref{eq:q2-prop}，几何与初值不变，仍不显式计入潜热和收缩。目标是确定使药材各处含水率低于 $0.15\,\mathrm{kg/kg}$ 的烘干时间，而不是另建优化模型。

附件环境序列只覆盖 $0$--$4\,\mathrm{h}$。末一小时温度与水分浓度的均值分别为 $50.00^\circ\mathrm{C}$ 与 $0.0500\,\mathrm{kg/kg}$，与题面所述恒温干燥阶段一致。因此 $0$--$4\,\mathrm{h}$ 仍使用观测点的 PCHIP 曲线，$4\,\mathrm{h}$ 之后取
\[
T_a(t)=50^\circ\mathrm{C},\qquad C_a(t)=0.05\,\mathrm{kg/kg}.
\]
切换处允许环境输入出现小幅跳变，药材内部 $T,C$ 保持连续。该延续是有数据支持的工作假设，不是未来实测。

$50^\circ\mathrm{C}$ 时，附录扩散系数在 $C=2.55$ 约为 $1.35\times10^{-8}\,\mathrm{m^2/s}$，在 $C=0.15$ 约为 $8.00\times10^{-10}\,\mathrm{m^2/s}$，后者只有前者的约 $6\%$。后期内部变干后扩散显著变慢，不能用初始 $D$ 或前 $3\,\mathrm{h}$ 的下降斜率外推总时长。

\subsection{达标判据}

定义径向最大含水率
\begin{equation}
M(t)=\max_{0\le r\le R}C(r,t),\qquad
t_*:=\inf\{t\ge0:M(t)<0.15\}.
\label{eq:q3-event}
\end{equation}
``各处低于阈值''必须用 $M(t)$ 判断。表格中的五个半径以及结果文件中的 $21$ 个采样点只是输出位置，不能代替全域最大值。若解沿径向向外不增，则 $M(t)=C(0,t)$，但仍需在加密网格上检查最大位置。

连续向下穿越时，临界时刻满足 $M(t_*)=0.15$；题目要求的是严格小于。正式报告时刻 $t_{\mathrm{rep}}$ 取略晚于 $t_*$ 且经未舍入场复核满足 $M(t_{\mathrm{rep}})<0.15$ 的时刻，并检查数值误差，不能只把临界根四舍五入后宣称已经达标。题面中的 $2$--$3$ 天是工艺经验范围，不作为强制截断或拟合目标。

\subsection{求解与输出}

数值格式与问题二相同：表面加密的径向有限体积、物性随局部状态更新、联合 BDF。在阈值附近限制步长，用符号变化括号定位事件，避免把输出间隔 $60\,\mathrm{s}$ 误当成求解步长。温度场作为内部状态保留，即使题目只要求导出含水率。

表~\ref{tab:q3C} 按题目格式列出每隔 $6\,\mathrm{h}$ 及结束时刻的含水率。当前全文稿先给出模型、判据与环境假设；结束时刻的具体小时数及表中数值，待按上述方程完成全程积分后填入，不在此处估计或编造。

\begin{table}[htbp]
\centering
\caption{药材烘干过程的水分浓度（单位：$\mathrm{kg/kg}$，结束时刻待补）}
\label{tab:q3C}
\begin{tabular}{cccccc}
\toprule
时间/h & $0\,\mathrm{cm}$ & $0.5\,\mathrm{cm}$ & $1.0\,\mathrm{cm}$ & $1.5\,\mathrm{cm}$ & $2.0\,\mathrm{cm}$ \\
\midrule
6  & --- & --- & --- & --- & --- \\
12 & --- & --- & --- & --- & --- \\
18 & --- & --- & --- & --- & --- \\
$\cdots$ & --- & --- & --- & --- & --- \\
结束时刻 & --- & --- & --- & --- & --- \\
\bottomrule
\end{tabular}
\end{table}
